\begin{center}
\textbf{Accélération d'un véhicule}
\end{center}

Un véhicule électrique se déplace sur une route horizontale et rectiligne.

\begin{center}
\begin{pspicture}(2,-0.5)(10,0.5)
\psset{arrowsize=0.3}
\psline[linewidth=1.2pt]{->}(10,0)(2,0)
\uput[u](6,0){Sens du mouvement}
\end{pspicture}
%\includegraphics[scale=0.5]{voiture_polynesie_2025}

\psset{xunit=1cm,yunit=0.9cm}
\scalebox{0.7}{
\begin{pspicture}(0,1)(12,5)
%\psgrid
\pscircle[linewidth=8pt](9,1.8){0.8}
\pscircle[linewidth=8pt](3,1.8){0.8}
\psline[linewidth=4pt,linecolor=white](2.2,1.8)(9.8,1.8)
\pscircle[linewidth=1pt,fillstyle=solid,fillcolor=black](3,1.8){0.3}
\pscircle[linewidth=1pt,fillstyle=solid,fillcolor=black](9,1.8){0.3}
\psset{linewidth=3pt}
\psarc(9,1.8){1}{0}{180} \psline(8,1.8)(8,1.4) \psline(10,1.8)(10,1.6)
\psarc(3,1.8){1}{0}{180} \psline(8,1.8)(8,1.4) \psline(2,1.8)(2,1.4)
\psline(4,1.8)(4,1.4)(8,1.4) \psline(2,1.4)(1,1.6)(1,2)(1.8,2)
\pscurve(1,2.2)(1,2.4)(2,3)(4,3.4)(4,3.4)(6,4.4)(6,4.4)(7,4.6)(8,4.6)(10,3.8)(11,3)
\psline(11,3)(10.4,3) \psarc(11,3){0.6}{180}{300}
\psline(11,3)(11.4,2.2)(11.4,1.8)(10,1.6)
\pscurve(2,3)(1.5,2.3)(1,2.2)
\psline(4.4,1.4)(4.4,3.4)(9.6,3.4)
\pscurve(4.4,3.4)(6,4.2)(6,4.2)(7,4.4)(8,4.4)(9,4.1)(9.6,3.7)(9.6,3.4)(9.6,3.4)
\pscurve(7.6,3.4)(7.8,3)(7,1.4)
\psline(8,4.4)(7.6,3.4) \psline(7,2.8)(7.4,2.8) \psline(1.6,2.2)(11,2.2)
\end{pspicture}
}%%% fin du scalebox
\end{center}

Le graphique du \textbf{document \no 1 } représente l'évolution de la vitesse instantanée $v(t)$ (exprimée en $\text{m}\cdot\text{s}^{-1}$ en fonction du temps $t$ (exprimé en seconde) durant 50 secondes.

\textbf{1 u.a. = 1 unité d'aire.}

\begin{center}
\psset{xunit=0.2cm,yunit=0.2cm,labelFontSize=\scriptstyle,arrowsize=2pt 3}
\begin{pspicture}(-10,-5)(57,27)
\psgrid[unit=1cm,subgriddiv=5,gridlabels=0,gridcolor=gray,subgridcolor=lightgray](0,0)(11,5)
\psaxes[linewidth=0.95pt,Dx=5,Dy=5,]{->}(0,0)(55,25)
\psframe[fillstyle=hlines,hatchwidth=0.1pt=](45,15)(50,20)
\uput[d](55,0){\footnotesize $t$ en $s$}
\uput[l](-4,10){\rotatebox{90}{\small Vitesse en $\text{m}\cdot\text{s}^{-1}$}}
\uput[u](47.5,16){\footnotesize 25 u.a.}
\uput[d](25,-4){\textbf{Document \no 1 :} évolution de la vitesse instantanée}
\psset{linewidth=1.25pt,linecolor=blue}
\psbezier(0,0)(3,0)(2,10)(5,10) \psline(5,10)(10,10)
\psbezier(10,10)(17,10)(23,20)(30,20) \psbezier(30,20)(37,20)(43,0)(50,0)
\end{pspicture}
\end{center}

L'accélération instantanée est la dérivée de la vitesse par rapport au temps:
\[a(t)=\dfrac{\mathrm{d}v}{\mathrm{d}t}(t)\]
avec :
\begin{itemize}[label=$\bullet$]
\item $v(t)$ la vitesse instantanée en m$\cdot$s$^{-1}$ à l'instant $t$;
\item $a(t)$ l'accélération instantanée en m$\cdot$s$^{-2}$ à l'instant $t$.
\end{itemize}

On s'intéresse à l'accélération instantanée du véhicule sur l'intervalle [1 ; 4], et on admet que sur cet intervalle, on a : $v(t) = 2t$.

\begin{enumerate}[start=2]
\item Déterminer la valeur de l'accélération $a(t)$ sur l'intervalle [1 ; 4].
\end{enumerate}
La distance totale, notée $D$ et exprimée en mètre, parcourue par le véhicule en 50 secondes est donnée par :
\[D = \int_0^{50} v(t)\d t\]

\begin{enumerate}[resume]
\item Donner, en exploitant le graphique, une estimation de la distance totale parcourue par le véhicule.
\item Déduire la valeur de la vitesse moyenne du véhicule, exprimée en $\text{km}\cdot\text{h}^{-1}$, dans l'intervalle de temps [0 ; 50].
\end{enumerate}

\bigskip
